Actualmente, la organización opera con múltiples fuentes de datos distribuidas en diferentes sistemas, incluyendo el ERP, archivos Excel y otras plataformas operativas, lo que genera inconsistencias, duplicidad de información y procesos manuales en la consolidación de datos financieros. 

Esta situación provoca falta de confiabilidad en la información financiera, alto consumo de tiempo en tareas manuales de conciliación y validación de datos, retrasos en la generación de reportes financieros, mayor riesgo de errores humanos y dificultad para acceder a información consolidada en tiempo real. Además, la ausencia de una arquitectura de datos centralizada dificulta la automatización de procesos y limita la capacidad del área financiera para realizar análisis eficientes y oportunos.
